\section{Билет 37. Понятие изопроцесса. Газовые законы Бойля–Мариотта, Гей-Люссака и Шарля. Диаграммы каждого изопроцесса в осях $pV$, $VT$, $pT$.}

\begin{definition}
\textbf{Изопроцессом} называется термодинамический процесс, протекающий при постоянном значении одного из макроскопических параметров состояния системы (при неизменном количестве вещества $\nu = \text{const}$).
Для идеального газа, состояние которого описывается уравнением Менделеева--Клапейрона $pV = \nu RT$, выделяют три основных изопроцесса:
\begin{itemize}
    \item \textbf{Изотермический} ($T = \text{const}$);
    \item \textbf{Изобарный} ($p = \text{const}$);
    \item \textbf{Изохорный} ($V = \text{const}$).
\end{itemize}
\end{definition}

\begin{theorem}[Закон Бойля--Мариотта --- изотермический процесс]
При постоянной температуре $T$ и неизменном количестве вещества $\nu$ произведение давления газа на его объём остаётся постоянным:
\begin{equation}
    pV = \text{const} \quad (T = \text{const}). \label{eq:boyle_mariotte_37}
\end{equation}
\textbf{Графическое представление:}
\begin{itemize}
    \item В координатах $(p, V)$ --- гипербола ($p \sim 1/V$), называемая \textbf{изотермой}. Чем выше температура, тем выше расположена изотерма.
    \item В координатах $(V, T)$ --- вертикальная прямая ($T = \text{const}$).
    \item В координатах $(p, T)$ --- горизонтальная прямая ($T = \text{const}$).
\end{itemize}
\end{theorem}

\begin{theorem}[Закон Гей-Люссака --- изобарный процесс]
При постоянном давлении $p$ и неизменном количестве вещества $\nu$ отношение объёма газа к его абсолютной температуре остаётся постоянным:
\begin{equation}
    \frac{V}{T} = \text{const} \quad (p = \text{const}). \label{eq:gay_lussac_37}
\end{equation}
\textbf{Графическое представление:}
\begin{itemize}
    \item В координатах $(V, T)$ --- прямая, проходящая через начало координат ($V \sim T$), называемая \textbf{изобарой}. Чем больше давление, тем ниже расположена изобара.
    \item В координатах $(p, V)$ --- горизонтальная прямая ($p = \text{const}$).
    \item В координатах $(p, T)$ --- вертикальная прямая ($p = \text{const}$).
\end{itemize}
\end{theorem}

\begin{theorem}[Закон Шарля --- изохорный процесс]
При постоянном объёме $V$ и неизменном количестве вещества $\nu$ отношение давления газа к его абсолютной температуре остаётся постоянным:
\begin{equation}
    \frac{p}{T} = \text{const} \quad (V = \text{const}). \label{eq:charles_37}
\end{equation}
\textbf{Графическое представление:}
\begin{itemize}
    \item В координатах $(p, T)$ --- прямая, проходящая через начало координат ($p \sim T$), называемая \textbf{изохорой}. Чем больше объём, тем ниже расположена изохора.
    \item В координатах $(p, V)$ --- вертикальная прямая ($V = \text{const}$).
    \item В координатах $(V, T)$ --- горизонтальная прямая ($V = \text{const}$).
\end{itemize}
\end{theorem}

\begin{remark}
Все три закона справедливы для \textbf{идеального газа} при условии, что масса газа не изменяется. В реальных газах при высоких давлениях и низких температурах наблюдаются отклонения от этих законов из-за учёта собственного объёма молекул и сил межмолекулярного взаимодействия (уравнение Ван-дер-Ваальса).
\end{remark}


\begin{figure}[H]
\centering
\begin{tikzpicture}

% Изотерма
\begin{axis}[
width=5cm,
height=4cm,
at={(0,0)},
xlabel={$V$},
ylabel={$p$},
title={Изотерма},
xmin=0.5,xmax=5,
ymin=0,ymax=5,
ticks=none
]
\addplot[thick,domain=0.7:5,samples=100]
{3/x};
\end{axis}

% Изобара
\begin{axis}[
width=5cm,
height=4cm,
at={(6cm,0)},
xlabel={$T$},
ylabel={$V$},
title={Изобара},
xmin=0,xmax=5,
ymin=0,ymax=5,
ticks=none
]
\addplot[thick]
coordinates {(0,0) (4,4)};
\end{axis}

% Изохора
\begin{axis}[
width=5cm,
height=4cm,
at={(12cm,0)},
xlabel={$T$},
ylabel={$p$},
title={Изохора},
xmin=0,xmax=5,
ymin=0,ymax=5,
ticks=none
]
\addplot[thick]
coordinates {(0,0) (4,4)};
\end{axis}

\end{tikzpicture}
\caption{Графики основных изопроцессов идеального газа.}
\end{figure}
